\subsection{Independent extension through $N=2^{14}$}
\label{subsec:numerics-16384}

A separate memory-limited computation extended the exact checks and structural diagnostics to
\begin{equation}
 N=2^{14}=16384,
 \qquad
 X_N=(16385)^2-1\approx2.68\times10^8.
 \label{eq:extended-range-16384}
\end{equation}
The larger targets $N=2^{16}$ and $2^{17}$ would require sieving past approximately $4.3\times10^9$ and $1.7\times10^{10}$, respectively, and were not attempted in the available memory.  No values at those scales are claimed.

The independently constructed smooth, transport, and residual values agree exactly at every tested checkpoint.  At the terminal checkpoint,
\begin{equation}
 \boxed{
 A_{16384}=-851896,
 \qquad
 T_{16384}=-850299,
 \qquad
 S_{16384}=-1597=A_{16384}-T_{16384}.
 }
 \label{eq:extended-terminal-identity}
\end{equation}

For a cofactor block $C<c\le2C$ with $C\asymp N^\sigma$, define the theorem-predicted rank-one residual
\begin{equation}
 E_C^{\mathrm{th}}
 :=B_C-\lambda_\sigma a_C P,
 \qquad
 \lambda_\sigma=\frac{2}{2-\sigma}.
 \label{eq:theorem-predicted-residual}
\end{equation}
This is distinct from the orthogonal least-squares residual $B_C-\widehat\beta_CP$.  The following table reports the former.

\begin{table}[htbp]
\centering
\caption{Theorem-predicted rank-one residuals at $N=16384$.  Residual energy is four to five orders of magnitude below the raw block energy throughout the tested range.}
\label{tab:extended-rank-one-residuals}
\scriptsize
\setlength{\tabcolsep}{4.5pt}
\begin{tabular}{lrrrr}
\toprule
cofactor block & $a_C$ & $\|B_C\|^2/N^3$ & $\|E_C^{\mathrm{th}}\|^2/N^3$ & residual/raw\\
\midrule
$(1,2]$     & $-0.50000$ & 45202.3 & 0.5166 & $1.1\times10^{-5}$\\
$(2,4]$     & $-0.33333$ & 21074.3 & 2.4454 & $1.2\times10^{-4}$\\
$(4,8]$     & $-0.17619$ &  6344.8 & 1.0401 & $1.6\times10^{-4}$\\
$(8,16]$    & $+0.07026$ &  1133.5 & 0.0481 & $4.2\times10^{-5}$\\
$(16,32]$   & $-0.12347$ &  3752.6 & 0.1394 & $3.7\times10^{-5}$\\
$(32,64]$   & $+0.07872$ &  1626.2 & 0.4517 & $2.8\times10^{-4}$\\
$(64,128]$  & $-0.01655$ &    77.4 & 0.0547 & $7.1\times10^{-4}$\\
$(128,256]$ & $+0.00512$ &     8.9 & 0.0052 & $5.8\times10^{-4}$\\
$(256,512]$ & $-0.00911$ &    29.4 & 0.0058 & $2.0\times10^{-4}$\\
\bottomrule
\end{tabular}
\end{table}

The height-shell decomposition is more decisive.  With shells defined by the ratio $|Y_m|/j$, where $j$ is the entry block, the diagonal shell energies are:

\begin{table}[htbp]
\centering
\caption{Height-shell diagonal energies at $N=16384$.  The recombined energy is not the sum of these rows; it includes their large signed cross terms.}
\label{tab:extended-height-shells}
\small
\begin{tabular}{lr}
\toprule
shell for $|Y_m|/j$ & $\sum_n|S_n^{[\ell]}|^2/N^3$\\
\midrule
$[0,5)$       & 0.00057\\
$[5,20)$      & 0.0147\\
$[20,80)$     & 0.259\\
$[80,320)$    & 3.81\\
$[320,1280)$  & 44.7\\
$[1280,\infty)$ & 86.3\\
\midrule
signed recombination & 0.0104\\
\bottomrule
\end{tabular}
\end{table}

Using the rounded diagonal values in \cref{tab:extended-height-shells},
\begin{equation}
 \sum_\ell G_{\ell\ell}(N)/N^3\approx135.084,
 \qquad
 2\Repart\sum_{\ell<m}G_{\ell m}(N)/N^3\approx-135.074.
 \label{eq:extended-cross-shell-cancellation}
\end{equation}
Thus the high sector is not small by shell.  Its observed smallness is produced by nearly complete cancellation of large off-diagonal shell interactions.  Any proof that takes absolute values or sums separate shell norms before recombination destroys the principal mechanism visible in the data.

\subsubsection{Exact and secondary reduced-square modes}

The exact congruence $q^2\equiv1\pmod{24}$ was verified on all $4459$ primes in the tested prime sample.  Modulo $40$, the same sample occupied exactly the two classes
\begin{equation}
 q^2\equiv1\text{ or }9\pmod{40},
 \label{eq:prime-square-mod40}
\end{equation}
with counts $2224$ and $2235$.  This is forced for primes $q>5$ by $q^2\equiv1\pmod8$ and $q^2\equiv1$ or $4\pmod5$.

\begin{proposition}[Two reduced-square classes modulo $40$]
\label{prop:prime-square-mod40}
For every prime $q>5$,
\begin{equation}
 \boxed{q^2\equiv1\text{ or }9\pmod{40}.}
 \label{eq:prime-square-mod40-boxed}
\end{equation}
\end{proposition}

Accordingly, a conductor dividing $40$ produces a two-mode prime phase rather than a featureless minor-arc sum.  In the tested sample,
\begin{center}
\begin{tabular}{lc}
\toprule
frequency & observed prime-sum magnitude\\
\midrule
$1/20$ & $3607$\\
$3/20$ & $1378$\\
$1/40$ & approximately $8$--$14$\\
$9/40$ & approximately $8$--$14$\\
\bottomrule
\end{tabular}
\end{center}
The same two exact phase modes can therefore interfere constructively or destructively.  Here ``resonant'' means low-rank exact phase support, not necessarily large amplitude.

Cumulative M\"obius-weighted cofactor phase magnitudes through the six blocks ending at $c=64$ were
\begin{equation}
\begin{array}{c|rrrrrr}
\theta & (1,2]&(2,4]&(4,8]&(8,16]&(16,32]&(32,64]\\
\hline
1/12 & 1.0&1.6&3.4&4.4&8.0&13.9\\
1/20 & 1.0&1.8&2.5&5.2&7.3&10.7\\
1/40 & 1.0&2.0&1.8&2.2&0.38&2.8
\end{array}
\label{eq:extended-cofactor-phase-cascade}
\end{equation}
The $1/40$ row contains a substantial cancellation event at the block $16<c\le32$, emphasizing again that the resonant main term is a signed cofactor cascade rather than a monotone positive contribution.
